% ============================================================
%  CSE200 – Technical Writing and Presentation
%  BUET | Department of CSE
%  Report Skeleton — matches senior batch Autoencoder style
% ============================================================

\documentclass[12pt, a4paper]{article}

% ── Packages ─────────────────────────────────────────────────
\usepackage[T1]{fontenc}
\usepackage[utf8]{inputenc}
\usepackage{lmodern}                % Clean Computer Modern font (matches senior batch)
\usepackage[margin=1in]{geometry}   % Standard 1-inch margins
\usepackage{graphicx}               % \includegraphics
\usepackage{amsmath, amssymb}       % Math equations
\usepackage{booktabs}               % Professional tables (\toprule etc.)
\usepackage{array}                  % Better column control in tables
\usepackage{hyperref}               % Blue hyperlinks (TOC, refs)
\usepackage{xcolor}                 % Color support
\usepackage{float}                  % [H] figure placement
\usepackage{caption}                % Custom caption formatting
\usepackage{subcaption}             % Side-by-side figures
\usepackage{parskip}                % Space between paragraphs, no indent
\usepackage{setspace}               % Line spacing
\usepackage{titlesec}               % Section heading style
\usepackage{enumitem}               % Custom lists
\usepackage{fancyhdr}               % Header/footer
\usepackage{lipsum}                 % Placeholder text — REMOVE when writing

% ── Hyperlink colors (matches the blue TOC in senior batch) ──
\hypersetup{
    colorlinks = true,
    linkcolor  = {blue!60!black},
    citecolor  = {blue!60!black},
    urlcolor   = {blue!60!black}
}

% ── Caption style ────────────────────────────────────────────
\captionsetup{
    font       = small,
    labelfont  = bf,
    format     = plain,
    justification = centering
}

% ── Section heading style ────────────────────────────────────
\titleformat{\section}{\large\bfseries}{\thesection}{1em}{}
\titleformat{\subsection}{\normalsize\bfseries}{\thesubsection}{1em}{}
\titleformat{\subsubsection}{\normalsize\itshape}{\thesubsubsection}{1em}{}

% ── Page style ───────────────────────────────────────────────
\pagestyle{fancy}
\fancyhf{}
\renewcommand{\headrulewidth}{0pt}
\cfoot{\thepage}

% ── Paragraph spacing ────────────────────────────────────────
\setlength{\parskip}{6pt}
\setlength{\parindent}{0pt}

% ============================================================
%  DOCUMENT START
% ============================================================
\begin{document}

% ── COVER PAGE ───────────────────────────────────────────────
\begin{titlepage}
    \centering
    \vspace*{0.5cm}

    % University logo — replace buet_logo.png with your actual file
    \includegraphics[width=0.35\textwidth]{images/buet_logo.png}\\[1.2cm]

    {\large\textbf{CSE200 - Technical Writing and Presentation}}\\[0.8cm]

    {\Large\textbf{BANGLADESH UNIVERSITY OF\\[0.2cm]
    ENGINEERING AND TECHNOLOGY}}\\[0.4cm]

    {\normalsize\textbf{DEPARTMENT OF COMPUTER SCIENCE AND ENGINEERING}}\\[2.5cm]

    \rule{\textwidth}{1.5pt}\\[0.3cm]
    {\LARGE\textbf{YOUR TOPIC TITLE HERE}}\\[0.3cm]
    \rule{\textwidth}{1.5pt}\\[2cm]

    {\normalsize 23 February 2026}\\[1.5cm]

    % Group info table — left-aligned block centred on page
    \begin{minipage}{0.5\textwidth}
        \begin{tabbing}
            \hspace{3.5cm} \= \kill
            Subsection      \> : B1 \\
            Group           \> : 4 \\
            Md. Rashidul Islam   \> : 2205072 \\
            Fathum Mobin   \> : 2205073 \\
            Md. Riyadun Nabi   \> : 2205076 \\
        \end{tabbing}
    \end{minipage}

\end{titlepage}

% ── TABLE OF CONTENTS ────────────────────────────────────────
\tableofcontents
\thispagestyle{fancy}
\newpage

% ── ABSTRACT ─────────────────────────────────────────────────
\newpage
\begin{center}
    {\large\textbf{Abstract}}
\end{center}

% Write 150–200 words summarising: what the topic is, why it matters,
% key components/methods, notable variants, challenges, real-world impact.
\lipsum[1]   % ← Replace with your abstract

\newpage

% ============================================================
%  SECTION 1 — Introduction
% ============================================================
\section{Introduction}

% Motivate the topic with a real-world analogy. State what it does and
% why it was developed. End with a brief outline of the report.
\lipsum[2]   % ← Replace with your introduction

Key capabilities include:
\begin{itemize}
    \item \textbf{Capability 1}: One-line description.
    \item \textbf{Capability 2}: One-line description.
    \item \textbf{Capability 3}: One-line description.
\end{itemize}

% ============================================================
%  SECTION 2 — Background / Overview
% ============================================================
\section{Background and Overview}

\lipsum[3]   % ← Replace with background

\subsection{Sub-topic A}
\lipsum[4]

\subsection{Sub-topic B}
\lipsum[5]

% ============================================================
%  SECTION 3 — Core Architecture / Concepts
% ============================================================
\section{Core Architecture}

\lipsum[6]

% ── Figure placeholder ────────────────────────────────────────
\begin{figure}[H]
    \centering
    % \includegraphics[width=0.8\textwidth]{architecture.png}
    \fbox{\rule{0pt}{5cm}\rule{0.8\textwidth}{0pt}}  % ← Remove this line when you add real figure
    \caption{Overall architecture of \textit{Your Topic}. The diagram shows [describe what it shows].}
    \label{fig:architecture}
\end{figure}

\subsection{Component 1}
\lipsum[7]

\subsection{Component 2}
\lipsum[8]

% ── Sub-figure example (side-by-side) ────────────────────────
\begin{figure}[H]
    \centering
    \begin{subfigure}[b]{0.45\textwidth}
        \centering
        % \includegraphics[width=\textwidth]{fig_a.png}
        \fbox{\rule{0pt}{3cm}\rule{\textwidth}{0pt}}
        \caption{Left subfigure caption.}
        \label{fig:a}
    \end{subfigure}
    \hfill
    \begin{subfigure}[b]{0.45\textwidth}
        \centering
        % \includegraphics[width=\textwidth]{fig_b.png}
        \fbox{\rule{0pt}{3cm}\rule{\textwidth}{0pt}}
        \caption{Right subfigure caption.}
        \label{fig:b}
    \end{subfigure}
    \caption{Comparison of (a) [left] and (b) [right] showing [describe difference].}
    \label{fig:comparison}
\end{figure}

\subsection{Component 3}
\lipsum[9]

% ============================================================
%  SECTION 4 — Mathematical Formulation
% ============================================================
\section{Mathematical Formulation}

% Introduce all notation BEFORE presenting equations.
% Always number every equation and reference them inline as Eq.~\eqref{eq:loss}.

Let $\mathbf{x} \in \mathbb{R}^n$ denote the input. The primary objective is to minimise the loss function:

\begin{equation}
    \mathcal{L}(\theta) = \frac{1}{N} \sum_{i=1}^{N} \ell\!\left(y_i,\, \hat{y}_i\right)
    \label{eq:loss}
\end{equation}

where $\theta$ are the model parameters, $y_i$ is the ground truth, and $\hat{y}_i$ is the prediction.

The parameter update rule follows gradient descent:

\begin{equation}
    \theta \leftarrow \theta - \eta \, \nabla_\theta \mathcal{L}(\theta)
    \label{eq:update}
\end{equation}

where $\eta$ is the learning rate.

% ── Aligned multi-line derivation example ────────────────────
\begin{align}
    \log p(\mathbf{x})
        &= \log \int p(\mathbf{x}, \mathbf{z})\, d\mathbf{z} \notag \\
        &\geq \mathbb{E}_{q(\mathbf{z}|\mathbf{x})}
              \left[\log \frac{p(\mathbf{x}, \mathbf{z})}{q(\mathbf{z}|\mathbf{x})}\right]
        \label{eq:elbo}
\end{align}

% ── Notation table ────────────────────────────────────────────
\begin{table}[H]
    \centering
    \caption{Summary of mathematical notation used in this report.}
    \label{tab:notation}
    \begin{tabular}{@{}cl@{}}
        \toprule
        \textbf{Symbol} & \textbf{Description} \\
        \midrule
        $\mathbf{x}$    & Input vector \\
        $\mathbf{z}$    & Latent / hidden variable \\
        $\theta$        & Model parameters \\
        $\eta$          & Learning rate \\
        $\mathcal{L}$   & Loss function \\
        $N$             & Number of training samples \\
        \bottomrule
    \end{tabular}
\end{table}

% ============================================================
%  SECTION 5 — Comparative Analysis
% ============================================================
\section{Comparative Analysis}

\lipsum[10]

\subsection{Variant / Approach A}
\lipsum[11]

\subsection{Variant / Approach B}
\lipsum[12]

\subsection{Comparison}

% ── Comparison table ─────────────────────────────────────────
\begin{table}[H]
    \centering
    \caption{Comparison of different approaches across key properties.}
    \label{tab:comparison}
    \begin{tabular}{@{}lcccc@{}}
        \toprule
        \textbf{Method} & \textbf{Property A} & \textbf{Property B}
                        & \textbf{Property C} & \textbf{Best for} \\
        \midrule
        Method 1  & High   & Low    & Yes & Use case 1 \\
        Method 2  & Medium & Medium & No  & Use case 2 \\
        Method 3  & Low    & High   & Yes & Use case 3 \\
        \bottomrule
    \end{tabular}
\end{table}

% ============================================================
%  SECTION 6 — Applications and Use Cases
% ============================================================
\section{Applications and Use Cases}

\lipsum[13]

\subsection{Application Area 1}
\lipsum[14]

\subsection{Application Area 2}
\lipsum[15]

% ============================================================
%  SECTION 7 — Advantages and Limitations
% ============================================================
\section{Advantages and Limitations}

\subsection{Advantages}
\begin{itemize}
    \item \textbf{Advantage 1}: Explanation.
    \item \textbf{Advantage 2}: Explanation.
    \item \textbf{Advantage 3}: Explanation.
\end{itemize}

\subsection{Limitations}
\begin{itemize}
    \item \textbf{Limitation 1}: Explanation.
    \item \textbf{Limitation 2}: Explanation.
\end{itemize}

% ============================================================
%  SECTION 8 — Advanced Variants / Extensions
% ============================================================
\section{Advanced Variants and Extensions}

\lipsum[16]

\subsection{Variant 1}
\lipsum[17]

\subsection{Variant 2}
\lipsum[18]

% ============================================================
%  SECTION 9 — Conclusion
% ============================================================
\section{Conclusion}

% Summarise key findings. Reflect on limitations and future directions.
% End with a forward-looking statement.
\lipsum[19]

% ============================================================
%  REFERENCES
% ============================================================
\newpage
\begin{thebibliography}{99}

\bibitem{ref1}
Author, A., Author, B., and Author, C.
\textit{Title of the Paper}.
Journal/Conference Name, Volume(Issue), pp.\ 1--10, Year.

\bibitem{ref2}
Author, D.
\textit{Title of the Book}.
Publisher, City, Year.

\bibitem{ref3}
Author, E.
``Title of Online Article.''
Website Name. [Online].
Available: \url{https://www.example.com}.
Accessed: Month DD, Year.

\end{thebibliography}

\end{document}